\documentclass[11pt,a4paper]{article}
\usepackage[T1]{fontenc}
\usepackage[utf8]{inputenc}
\usepackage{lmodern}
\usepackage{amsmath,amssymb,amsthm}
\usepackage{booktabs}
\usepackage{xcolor}
\usepackage[margin=27mm]{geometry}
\usepackage[colorlinks=true,urlcolor=blue,citecolor=blue,linkcolor=blue]{hyperref}
\newtheorem{theorem}{Theorem}
\newtheorem{proposition}[theorem]{Proposition}
\newtheorem{lemma}[theorem]{Lemma}
\newtheorem{corollary}[theorem]{Corollary}
\theoremstyle{definition}
\newtheorem{definition}[theorem]{Definition}
\newtheorem{conjecture}[theorem]{Conjecture}
\theoremstyle{plain}
\DeclareMathOperator{\Tr}{Tr}
\newcommand{\ket}[1]{\lvert#1\rangle}
\newcommand{\bra}[1]{\langle#1\rvert}
\newcommand{\ketbra}[2]{\lvert#1\rangle\langle#2\rvert}
\newcommand{\norm}[1]{\lVert#1\rVert}
\newcommand{\C}{\mathbb{C}}
\newcommand{\NS}{N_{S}}
\newcommand{\G}{\Gamma}
\setlength{\emergencystretch}{2em}
\setlength{\parindent}{0pt}
\setlength{\parskip}{0.3em plus 0.1em minus 0.1em}
\title{Structured negativity and negativity:\\
an exact criterion for their equality}
\author{Fable \& Astra\\[0.35em]
\normalsize Research led by Jean-Philippe Bruneton\\[-0.1em]
\small Universit\'e Paris Cit\'e, France}
\date{19 November 2026}

\begin{document}
\maketitle
\begin{abstract}
The partial transpose of an entangled state of two particles can have negative
eigenvalues, and the negativity, the total weight of these eigenvalues, is the
most widely used computable measure of entanglement. Because the partial
transpose is not a physical operation, Kumari and Adhikari introduced the
structured negativity, a quantity built from a physically implementable
approximation of the partial transpose, and conjectured that for two
$d$-level systems it coincides with the negativity whenever the partial
transpose has exactly $d(d-1)/2$ negative eigenvalues. We show that the
structured negativity equals $d$ times the largest negative eigenvalue in
absolute value, so that the two quantities agree exactly when the sum of the
negative eigenvalues equals $d(d-1)/2$ times the most negative one. The
conjecture therefore holds for two qubits and fails for every $d\ge3$: the
pure two-qutrit state $(2\ket{00}+2\ket{11}+\ket{22})/3$ has three negative
eigenvalues but negativity $8/9$ and structured negativity $4/3$. A family of
states whose partial transpose has $(d-1)^2$ equal negative eigenvalues, built
from the symmetric subspace of $2(d-1)$ qubits, saturates the inequality
$N\le 2(1-1/d)\NS$ of the same authors in every dimension, and so disproves
their further claim that the structured negativity always dominates the
negativity.
\end{abstract}

\section{Introduction}
Entanglement between two quantum systems can be detected by a simple matrix
operation. Transposing the second factor of a density matrix maps every
separable state to another state, so a negative eigenvalue of the partial
transpose certifies entanglement \cite{Peres1996,Horodecki1996}. The
negativity, the sum of the absolute values of these negative eigenvalues,
turns this test into an entanglement monotone that is computable from the
spectrum alone \cite{VidalWerner2002}. The partial transpose is, however, a
positive but not completely positive map, so it cannot be applied to a
quantum system in the laboratory. The structural physical approximation of
Horodecki and Ekert repairs this by mixing the map with just enough white
noise to make it completely positive \cite{HorodeckiEkert2002}; the spectrum
of the resulting physical state still carries the information about the
negative eigenvalues, shifted and rescaled.

Kumari and Adhikari built an entanglement quantifier on this idea
\cite{KumariAdhikari2022}. For a state $\rho$ of two $d$-level systems, their
structured negativity $\NS(\rho)$ is read off from the smallest eigenvalue of
the structural physical approximation of $\rho^{T_B}$. They compared it with
the negativity $N(\rho)$, normalised so that both quantities equal one on the
maximally entangled state, and, from four families of examples, conjectured
that the two coincide whenever the number $q$ of negative eigenvalues of
$\rho^{T_B}$ equals $d(d-1)/2$, that the inequality $N<\NS$ is strict for
other values of $q$, and that $\NS\ge N$ always.

The simplest example already shows what goes wrong. For a pure two-qutrit
state with Schmidt coefficients $\sqrt{s_1},\sqrt{s_2},\sqrt{s_3}$, the
partial transpose has three negative eigenvalues $-\sqrt{s_is_j}$, $i<j$. The
negativity sums them, while the structured negativity, as we show, sees only
the most negative one. The two can agree only if the three products
$s_is_j$ coincide, that is, only for the maximally entangled state. The
number of negative eigenvalues is thus the wrong invariant: what matters is
their \emph{effective} number, the ratio of their sum to the largest one.

Our main result makes this precise. The structured negativity equals $d$
times the absolute value of the smallest eigenvalue of $\rho^{T_B}$
(Proposition~\ref{prop:ns}), and $N=\NS$ holds if and only if the sum of the
negative eigenvalues of $\rho^{T_B}$ equals $d(d-1)/2$ times the most negative
one (Theorem~\ref{thm:criterion}). Consequently the conjecture is true and
trivial for two qubits and false for every $d\ge3$, with exact pure-state
counterexamples (Corollary~\ref{cor:pure}). In the opposite direction, we
construct in every dimension states whose partial transpose has the maximal
number $(d-1)^2$ of negative eigenvalues, all equal
(Theorem~\ref{thm:saturation}). They saturate the inequality
$N\le2(1-1/d)\NS$ proved in \cite{KumariAdhikari2022} and give $N>\NS$ for
every $d\ge3$. An explicit two-qutrit state with four negative eigenvalues and
$N=\NS$ shows that equality also occurs away from $q=d(d-1)/2$.

\section{Related work}\label{sec:related}
The partial transpose criterion is due to Peres \cite{Peres1996}, and its
sufficiency for two qubits and for a qubit and a qutrit to the Horodeckis
\cite{Horodecki1996}. Vidal and Werner showed that the negativity is an
entanglement monotone \cite{VidalWerner2002}. The structural physical
approximation was introduced by Horodecki and Ekert as a way to implement
positive maps, and in particular the partial transpose, by physical
operations followed by spectrum estimation \cite{HorodeckiEkert2002}.

The structured negativity and the conjecture discussed here are in Kumari
and Adhikari \cite{KumariAdhikari2022}, whose Result-1 states
$\NS\ge\frac{d(d-1)}{2q}N$ and, with the bound $q\le(d-1)^2$, $N\le
2(1-\frac1d)\NS$. The conjecture is repeated without resolution in Chapter~6 of
Kumari's thesis \cite{Kumari2023}. A later paper by Aggarwal, Kumari and
Adhikari applies the structural physical approximation to the realignment
criterion \cite{Aggarwal2023}; it does not return to the conjecture.

Our construction concerns the number and the size of the negative eigenvalues
of the partial transpose. Rana proved that the partial transpose of an
$m\otimes n$ state has at most $(m-1)(n-1)$ negative eigenvalues and that its
spectrum lies in $[-\frac12,1]$, extending the two-qubit result of Sanpera,
Tarrach and Vidal, and gave low-dimensional examples
\cite{Rana2013}. Johnston showed that the bound $(m-1)(n-1)$ is attained in
all dimensions \cite{Johnston2013}; his construction does not control the
individual negative eigenvalues. The states of Theorem~\ref{thm:saturation}
attain the bound for $m=n=d$ with all negative eigenvalues equal, which is
what the saturation of Result-1 requires. Jiang, Qiu and Yu recently
introduced the Ky Fan $k$-negativities, the sums of the $k$ most negative
eigenvalues of the partial transpose, and derived spectral bounds for them
\cite{JiangQiuYu2026}; in their notation Proposition~\ref{prop:ns} reads
$\NS=d\,\mathcal{N}_1$. They do not discuss the structured negativity.

\section{Definitions}\label{sec:defs}
Let $\rho$ be a state on $\C^d\otimes\C^d$ with $d\ge2$, and write
$\rho^\G=\rho^{T_B}$ for its partial transpose on the second factor in the
computational basis. Following \cite{KumariAdhikari2022}, the negativity is
normalised as
\begin{equation}\label{eq:N}
N(\rho)=\frac{\norm{\rho^\G}_1-1}{d-1}
=\frac{2}{d-1}\sum_{\lambda_i<0}\lvert\lambda_i(\rho^\G)\rvert ,
\end{equation}
where the second form uses $\Tr\rho^\G=1$. The structural physical
approximation of the partial transpose sends $\rho$ to the state
\begin{equation}\label{eq:spa}
\tilde\rho=\frac{d}{d^3+1}\,I\otimes I+\frac{1}{d^3+1}\,\rho^\G ,
\end{equation}
and the structured negativity is
\begin{equation}\label{eq:NS}
\NS(\rho)=K\max\Bigl\{\frac{d}{d^3+1}-\lambda_{\min}(\tilde\rho),\,0\Bigr\},
\qquad K=d(d^3+1).
\end{equation}
Both quantities equal one on the maximally entangled state
$\ket{\Phi}=d^{-1/2}\sum_i\ket{ii}$.

When $\rho^\G$ is not positive semidefinite, list its strictly negative
eigenvalues with multiplicity as $-\nu_1\le-\nu_2\le\cdots\le-\nu_q<0$, so
that $q$ is their number and $\nu_1=\lvert\lambda_{\min}(\rho^\G)\rvert$. We
call
\begin{equation}\label{eq:r}
r(\rho)=\frac{\nu_1+\nu_2+\cdots+\nu_q}{\nu_1}\in[1,q]
\end{equation}
the effective number of negative eigenvalues. It equals $q$ if and only if
all negative eigenvalues are equal.

\begin{conjecture}[Kumari and Adhikari \cite{KumariAdhikari2022}]
\label{conj:KA}
For a state of $\C^d\otimes\C^d$, $N(\rho)=\NS(\rho)$ when $q=d(d-1)/2$.
The authors also state that the inequality $N<\NS$ is strict when
$q\ne d(d-1)/2$ (Section~IV.A) and that $\NS\ge N$ always (Conclusion).
\end{conjecture}

\section{The exact criterion}\label{sec:criterion}
\begin{proposition}\label{prop:ns}
For every state $\rho$ of $\C^d\otimes\C^d$,
$\NS(\rho)=d\,\max\{0,-\lambda_{\min}(\rho^\G)\}$.
\end{proposition}
\begin{proof}
By \eqref{eq:spa}, $\lambda_{\min}(\tilde\rho)=(d+\lambda_{\min}(\rho^\G))/(d^3+1)$,
so $\frac{d}{d^3+1}-\lambda_{\min}(\tilde\rho)=-\lambda_{\min}(\rho^\G)/(d^3+1)$.
Multiplying by $K=d(d^3+1)$ gives the claim.
\end{proof}

The structural physical approximation is thus an affine rescaling of the
spectrum, and the structured negativity retains only its lowest point.

\begin{theorem}\label{thm:criterion}
For every state $\rho$ of $\C^d\otimes\C^d$ whose partial transpose is not
positive semidefinite,
\begin{equation}\label{eq:ratio}
\frac{N(\rho)}{\NS(\rho)}=\frac{2\,r(\rho)}{d(d-1)} .
\end{equation}
In particular $N(\rho)=\NS(\rho)$ if and only if
$\nu_1+\cdots+\nu_q=\frac{d(d-1)}{2}\,\nu_1$. Three regimes follow. If
$q<\frac{d(d-1)}{2}$, then $N<\NS$. If $q=\frac{d(d-1)}{2}$, then $N\le\NS$,
with equality if and only if all negative eigenvalues of $\rho^\G$ are equal.
If $q>\frac{d(d-1)}{2}$, then $N\le\frac{2q}{d(d-1)}\NS$, and each of $N<\NS$,
$N=\NS$ and $N>\NS$ can occur (Section~\ref{sec:beyond}).
\end{theorem}
\begin{proof}
Proposition~\ref{prop:ns} gives $\NS=d\nu_1$ and \eqref{eq:N} gives
$N=\frac{2}{d-1}\sum_i\nu_i$; their quotient is \eqref{eq:ratio}. The
regimes follow from $1\le r\le q$, with $r=q$ exactly when the $\nu_i$ are
all equal. For $d=3$ and $q=4$ the three cases are realised by the explicit
states of Section~\ref{sec:beyond}: $\rho_\star$ has $N>\NS$, $\rho_=$ has
$N=\NS$, and $\rho_<$ has $N<\NS$.
\end{proof}

Result-1 of \cite{KumariAdhikari2022}, $\NS\ge\frac{d(d-1)}{2q}N$, is
exactly the inequality $r\le q$. The conjecture asserts that
$q=\frac{d(d-1)}{2}$ forces $r=q$, which confuses the number of negative
eigenvalues with their effective number.

\begin{corollary}[two qubits]\label{cor:qubits}
For $d=2$, $N(\rho)=\NS(\rho)=2\max\{0,-\lambda_{\min}(\rho^\G)\}$ for every
state.
\end{corollary}
\begin{proof}
The partial transpose of a two-qubit state has at most one negative eigenvalue
\cite{Rana2013}, so $q\le1$ and $r=1=d(d-1)/2$ whenever $q\ge1$.
\end{proof}

For $d=2$ the conjecture therefore holds and says nothing beyond $q\le1$,
which is what the two-qubit examples of \cite{KumariAdhikari2022} (Werner
states and maximally entangled mixed states) display.

\begin{corollary}[pure states]\label{cor:pure}
Let $\ket\psi=\sum_{i=1}^d\sqrt{s_i}\,\ket{ii}$ with $s_i>0$ and
$\sum_is_i=1$. Then $q=\frac{d(d-1)}{2}$ and
\begin{equation}\label{eq:pure}
N=\frac{2}{d-1}\sum_{i<j}\sqrt{s_is_j},\qquad
\NS=d\,\max_{i<j}\sqrt{s_is_j}.
\end{equation}
For $d\ge3$, $N=\NS$ if and only if $\ket\psi$ is maximally entangled. The
conjecture is therefore false for every $d\ge3$.
\end{corollary}
\begin{proof}
The partial transpose
$\ketbra\psi\psi^\G=\sum_{i,j}\sqrt{s_is_j}\,\ketbra{ij}{ji}$ has eigenvalue
$s_i$ on $\ket{ii}$ and eigenvalues $\pm\sqrt{s_is_j}$ on
$(\ket{ij}\pm\ket{ji})/\sqrt2$ for $i<j$. This gives $q=d(d-1)/2$ and
\eqref{eq:pure}. By Theorem~\ref{thm:criterion}, equality holds if and only if
all products $s_is_j$, $i<j$, are equal. For $d\ge3$, comparing $s_is_k$ with
$s_js_k$ for a third index $k$ gives $s_i=s_j$ for all $i,j$.
\end{proof}

The simplest exact counterexample is
\begin{equation}\label{eq:psi}
\ket\psi=\tfrac13\bigl(2\ket{00}+2\ket{11}+\ket{22}\bigr)\in\C^3\otimes\C^3,
\end{equation}
whose partial transpose has negative eigenvalues $-\frac49,-\frac29,-\frac29$.
Here $q=3=\frac{d(d-1)}2$, $r=2$, $N=\frac89\approx0.889$ and
$\NS=\frac43\approx1.333$.

The conjecture is correct on families whose negative spectrum is flat with
$q=d(d-1)/2$. The isotropic states $x\ketbra\Phi\Phi+(1-x)I/d^2$ have a partial
transpose that is a combination of the identity and the swap operator, whose
only negative eigenvalue lives on the antisymmetric subspace of dimension
$d(d-1)/2$; the two-qutrit family $\rho_\alpha$ of \cite{KumariAdhikari2022}
has a partial transpose made of three identical $2\times2$ blocks. At the
other extreme, the Werner states have $q=1$ and $N/\NS=\frac{2}{d(d-1)}$.
These are the families from which the conjecture was inferred.

\section{Beyond \texorpdfstring{$d(d-1)/2$}{d(d-1)/2} negative eigenvalues}\label{sec:beyond}
\subsection{A flat negative spectrum of maximal size}
Set $n=d-1$ and identify $\C^d$ with the symmetric subspace
$\mathrm{Sym}^n(\C^2)$ of $n$ qubits through the Dicke basis
$\ket{k}=\binom nk^{-1/2}\sum_{\lvert x\rvert=k}\ket{x}$, $k=0,\dots,n$, where
$x$ runs over binary words of length $n$ with $k$ ones. Then
$\C^d\otimes\C^d\subset(\C^2)^{\otimes2n}$ contains $\mathrm{Sym}^{2n}(\C^2)$, of
dimension $2n+1=2d-1$, with orthonormal basis
\begin{equation}\label{eq:Dm}
\ket{D_m}=\sum_{k+l=m}\sqrt{\binom nk\binom nl\Big/\binom{2n}{m}}\;\ket{k}\ket{l},
\qquad m=0,\dots,2n .
\end{equation}
Let $Q$ be the orthogonal projector of $\C^d\otimes\C^d$ onto this subspace and
$P=I-Q$, of rank $d^2-(2d-1)=(d-1)^2$. Partial transposes are taken in the
Dicke basis.

\begin{lemma}\label{lem:QG}
The eigenvalues of $Q^\G$ are
\begin{equation}\label{eq:qJ}
q_J=\frac{(2n+1)\,(n!)^2}{(n-J)!\,(n+J+1)!},\qquad J=0,1,\dots,n,
\end{equation}
with multiplicity $2J+1$. They decrease strictly with $J$, and the smallest is
$q_n=\binom{2n}{n}^{-1}=\binom{2d-2}{d-1}^{-1}>0$.
\end{lemma}
\begin{proof}
For a unit vector $\psi\in\C^2$ write $\ket{\psi_n}=\psi^{\otimes n}\in
\mathrm{Sym}^n(\C^2)$, and let $d\psi$ be the uniform probability measure on
the unit sphere of $\C^2$. The operator
$(2n+1)\int(\ketbra\psi\psi)^{\otimes2n}d\psi$ is supported on
$\mathrm{Sym}^{2n}(\C^2)$, commutes with $U^{\otimes2n}$ for every
$U\in SU(2)$ and has trace $2n+1$. Since $\mathrm{Sym}^{2n}(\C^2)$ is an
irreducible representation, Schur's lemma identifies it with $Q$:
$Q=(2n+1)\int\ketbra{\psi_n}{\psi_n}\otimes\ketbra{\psi_n}{\psi_n}\,d\psi$. The
Dicke components of $\ket{\psi_n}$ are $\sqrt{\binom nk}\,\psi_0^{n-k}\psi_1^k$,
so complex conjugation in the Dicke basis sends $\ket{\psi_n}$ to
$\ket{\bar\psi_n}$, and
\[
Q^\G=(2n+1)\int\ketbra{\psi_n}{\psi_n}\otimes\ketbra{\bar\psi_n}{\bar\psi_n}\,d\psi .
\]
Identify $\C^d\otimes\C^d$ with the operators on $\C^d$ through
$\ket k\otimes\ket l\leftrightarrow\ketbra kl$, an isometry for the
Hilbert--Schmidt inner product. The vector $\ket\phi\otimes\ket{\bar\phi}$
corresponds to $\ketbra\phi\phi$, and its inner product with the vector of an
operator $A$ is $\bra\phi A\ket\phi$. Hence $Q^\G$ becomes the map
$\mathfrak B(A)=(2n+1)\int\bra{\psi_n}A\ket{\psi_n}\,\ketbra{\psi_n}{\psi_n}\,d\psi$.
It is self-adjoint and commutes with the adjoint action $A\mapsto U_nAU_n^\dagger$
of $SU(2)$, where $U_n$ is the restriction of $U^{\otimes n}$ to
$\mathrm{Sym}^n(\C^2)$. Under this action the operator space decomposes into
irreducible components of spin $J=0,\dots,n$, each of multiplicity one and
dimension $2J+1$ (the Clebsch--Gordan decomposition of two spins $n/2$). By
Schur's lemma $\mathfrak B$ acts on the spin-$J$ component as a scalar $q_J$,
and for every nonzero $A$ in that component
\[
q_J=\frac{(2n+1)\int\lvert\bra{\psi_n}A\ket{\psi_n}\rvert^2\,d\psi}{\Tr(A^\dagger A)} .
\]
Take $A=S_+^J$ with $S_+=\sum_{i=1}^n\sigma_+^{(i)}$ and
$\sigma_+=\ketbra01$. It is nonzero on $\mathrm{Sym}^n(\C^2)$ for $J\le n$,
has weight $J$ under the adjoint action ($[S_z,S_+^J]=JS_+^J$) and is
annihilated by the raising operator ($[S_+,S_+^J]=0$), so it is a highest
weight vector of spin $J$. Expanding $S_+^J$ as a sum over $J$-tuples of
sites, tuples with a repetition vanish because $\sigma_+^2=0$, and each of the
$n!/(n-J)!$ tuples of distinct sites contributes
$\bra\psi\sigma_+\ket\psi^J$; hence
$\bra{\psi_n}S_+^J\ket{\psi_n}=\frac{n!}{(n-J)!}(\bar\psi_0\psi_1)^J$. Since
$\lvert\psi_0\rvert^2$ is uniformly distributed on $[0,1]$,
$\int\lvert\psi_0\rvert^{2J}\lvert\psi_1\rvert^{2J}d\psi=\int_0^1x^J(1-x)^Jdx
=\frac{(J!)^2}{(2J+1)!}$. For the norm,
$S_+\ket k=\sqrt{k(n-k+1)}\,\ket{k-1}$ gives
$\norm{S_+^J\ket k}^2=\frac{k!}{(k-J)!}\frac{(n-k+J)!}{(n-k)!}
=(J!)^2\binom kJ\binom{n+J-k}{J}$, and the identity
$\sum_k\binom kJ\binom{M-k}{J}=\binom{M+1}{2J+1}$ with $M=n+J$ gives
$\Tr(S_-^JS_+^J)=(J!)^2\binom{n+J+1}{2J+1}$. Therefore
\[
q_J=(2n+1)\Bigl(\frac{n!}{(n-J)!}\Bigr)^2\frac{(J!)^2}{(2J+1)!}\cdot
\frac{(2J+1)!\,(n-J)!}{(J!)^2\,(n+J+1)!}
=\frac{(2n+1)(n!)^2}{(n-J)!\,(n+J+1)!}.
\]
Finally $q_J/q_{J+1}=\frac{n+J+2}{n-J}>1$ and
$q_n=\frac{(2n+1)(n!)^2}{(2n+1)!}=\binom{2n}{n}^{-1}$.
\end{proof}

\begin{theorem}\label{thm:saturation}
Let $d\ge2$, $\kappa_d=\binom{2d-2}{d-1}^{-1}$ and $1<t\le\frac{1}{1-\kappa_d}$.
The operator
\begin{equation}\label{eq:rhot}
\rho_t=\frac{I-t\,P^\G}{d^2-t(d-1)^2}
\end{equation}
is a state of $\C^d\otimes\C^d$, and its partial transpose
$\rho_t^\G=(I-tP)/(d^2-t(d-1)^2)$ has exactly $(d-1)^2$ negative eigenvalues,
all equal to $-\mu_t$ with $\mu_t=\frac{t-1}{d^2-t(d-1)^2}$. Consequently
\[
N(\rho_t)=2(d-1)\,\mu_t,\qquad \NS(\rho_t)=d\,\mu_t,\qquad
\frac{N(\rho_t)}{\NS(\rho_t)}=2\Bigl(1-\frac1d\Bigr).
\]
The inequality $N\le2(1-\frac1d)\NS$ of \cite{KumariAdhikari2022} is
saturated in every dimension, and for every $d\ge3$ there are states with
$N>\NS$.
\end{theorem}
\begin{proof}
Since $P^\G=I-Q^\G$, we have $I-tP^\G=tQ^\G-(t-1)I$, whose eigenvalues
$tq_J-(t-1)$ are at least $t\kappa_d-(t-1)\ge0$ by Lemma~\ref{lem:QG} and the
assumption $t(1-\kappa_d)\le1$. This operator is positive semidefinite and
nonzero, because $q_0=\frac{2n+1}{n+1}>1$; its trace $d^2-t(d-1)^2$ is
therefore positive and $\rho_t$ is a state. The partial transpose is an
involution, so $\rho_t^\G$ is proportional to $I-tP$, which equals $1$ on the
range of $Q$ and $1-t<0$ on the range of $P$, of dimension $(d-1)^2$. The values
of $N$ and $\NS$ follow from \eqref{eq:N} and Proposition~\ref{prop:ns}.
\end{proof}

At the endpoint $t=1/(1-\kappa_d)$ one finds
$\mu=\bigl((2d-1)\binom{2d-2}{d-1}-d^2\bigr)^{-1}$. For $d=2$ this state is
the Bell state $\ket{\Phi^+}$, with $N=\NS=1$. For $d=3$ ($t=\frac65$,
$\mu=\frac1{21}$) it reads, in the product basis $\ket{ij}$ of two qutrits,
\begin{multline}\label{eq:rhostar}
\rho_\star=\frac1{21}\Bigl[5\bigl(\ketbra{00}{00}+\ketbra{22}{22}\bigr)
+3\ketbra{11}{11}
+3\bigl(\ketbra{00}{11}+\ketbra{11}{22}+\mathrm{h.c.}\bigr)\\
+\bigl(\ketbra{00}{22}+\mathrm{h.c.}\bigr)
+2\ketbra uu+2\ketbra vv\Bigr],
\end{multline}
with $\ket u=\ket{10}+\ket{21}$ and $\ket v=\ket{01}+\ket{12}$. This rank-four
state has spectrum $\{\frac37,\frac4{21},\frac4{21},\frac4{21},0,0,0,0,0\}$ and
its partial transpose has spectrum $\{\frac5{21}\,(\times5),\,-\frac1{21}\,(\times4)\}$.
Hence $q=4$ and
\[
N(\rho_\star)=\frac4{21}\approx0.190\;>\;\NS(\rho_\star)=\frac17\approx0.143,
\]
contradicting the statement of \cite{KumariAdhikari2022} that the structured
negativity always dominates the negativity.

\subsection{Equality with \texorpdfstring{$q\ne d(d-1)/2$}{q != d(d-1)/2}}
The two-qutrit state
\begin{multline}\label{eq:rhoeq}
\rho_{=}=\frac1{64}\Bigl[21\bigl(\ketbra{00}{00}+\ketbra{22}{22}\bigr)
+2\ketbra{11}{11}
+5\bigl(\ketbra{00}{11}+\ketbra{11}{22}+\mathrm{h.c.}\bigr)
+4\bigl(\ketbra{00}{22}+\mathrm{h.c.}\bigr)\\
+4\ketbra uu+4\ketbra vv+2\ketbra{02}{02}+2\ketbra{20}{20}\Bigr]
\end{multline}
is positive semidefinite, with spectrum
$\{\frac{27}{64},\frac{17}{64},\frac18,\frac18,\frac1{32},\frac1{32},0,0,0\}$, and
its partial transpose has negative eigenvalues
$-\frac1{32},-\frac1{32},-\frac1{64},-\frac1{64}$. Thus $q=4\ne3$, $r=3$ and
$N=\NS=\frac3{32}$, so the claim that the inequality is strict whenever
$q\ne\frac{d(d-1)}{2}$ fails as well. The convex combination
$\rho_<=\frac1{10}\ketbra\psi\psi+\frac9{10}\rho_\star$ of the states
\eqref{eq:psi} and \eqref{eq:rhostar} also has $q=4$, with negative
eigenvalues $-\frac{11}{126}$, $-\frac{41}{630}$ (twice) and
$\frac5{42}-\frac{\sqrt{1747}}{315}\approx-0.0136$; there $r\approx2.65<3$
and $N\approx0.2311<\NS=\frac{11}{42}\approx0.2619$. Both $\rho_\star$ and $\rho_=$ are
invariant under $U\otimes\bar U$ for diagonal unitaries $U$. Their partial
transposes therefore split according to the value of $i+j$ into blocks of
sizes $1,2,3,2,1$, which were chosen of the form $-\varepsilon I+\ketbra ww$.

In summary, the conjecture ``$q=\frac{d(d-1)}2\Rightarrow N=\NS$'' holds for
$d=2$ and fails for every $d\ge3$; the correct statement is
Theorem~\ref{thm:criterion}, and its converse fails too. The inequality
$N\le\NS$ holds whenever $r\le\frac{d(d-1)}2$, in particular whenever
$q\le\frac{d(d-1)}2$, and is violated by $\rho_t$ for every $d\ge3$, where
$N/\NS$ reaches its largest possible value $2(1-\frac1d)$.

\section{Verification}\label{sec:verification}
The statements above are proved analytically; the counterexamples
\eqref{eq:psi}, \eqref{eq:rhostar} and \eqref{eq:rhoeq} have rational entries
and are checked in exact arithmetic. The public
\href{https://github.com/jpbruneton/open_quantum_problems/blob/main/research/entries/structured-negativity/checks/check_structured_negativity.py}{Python checker}
and its
\href{https://github.com/jpbruneton/open_quantum_problems/blob/main/research/entries/structured-negativity/data/verification.json}{recorded output}
form part of this entry. The checker computes $\NS$ literally from
definition \eqref{eq:NS}, through the state $\tilde\rho$, and tests
Proposition~\ref{prop:ns} and Theorem~\ref{thm:criterion} on 1149 random
states of ranks $1$ to $d^2$ with nonpositive partial transpose for
$d=2,3,4$, with largest deviation below $10^{-12}$; the equality $N=\NS$ for
two qubits; the Werner and isotropic families for $d=2,\dots,5$; the pure
counterexample \eqref{eq:psi} in exact arithmetic; 600 random pure states of
full Schmidt rank for $d=3,4,5$, all with $q=d(d-1)/2$ and $N<\NS$; 300 random
two-qutrit mixed states with exactly $q=3$, all with $N<\NS$; the exact
spectra of $\rho_\star$ and $\rho_=$ with SymPy; and, for $d=2,\dots,6$, the
spectrum of $Q^\G$ with multiplicities (projector built exactly, eigenvalues
compared in double precision with \eqref{eq:qJ}), the exact identity
$\rho_t^\G\propto I-tP$, the positivity of $\rho_t$ at $t=1/(1-\kappa_d)$ and
the ratio $N/\NS=2(d-1)/d$. The random samples only test conventions, such as
the normalisation of $N$ and the definition of $\NS$ through $\tilde\rho$;
they add no information beyond Theorem~\ref{thm:criterion}.

\section{Scope and limitations}\label{sec:scope}
Proposition~\ref{prop:ns}, Theorem~\ref{thm:criterion}, its corollaries,
Lemma~\ref{lem:QG} and Theorem~\ref{thm:saturation} are proved completely.
Theorem~\ref{thm:criterion} is an immediate rewriting of the definitions, and
the spectrum of the partial transpose of a pure state is classical. The
content of the note is the identification of the effective number $r(\rho)$
as the relevant invariant, the resulting refutation of the conjecture for
every $d\ge3$, and the family $\rho_t$.

The note leaves open the full set of values of $r(\rho)$ attainable at fixed
$q$. Theorem~\ref{thm:saturation} shows that $r=q$ is attained for
$q=(d-1)^2$ and the isotropic states attain it for $q=d(d-1)/2$; whether a
flat negative spectrum exists for intermediate values of $q$ has not been
examined. Nor do we describe the set of states with $r=d(d-1)/2$ beyond the
examples given. Other statements of \cite{KumariAdhikari2022}, in particular
the monotonicity of $\NS$ under local operations and classical communication,
are outside the scope of this note.

On bibliographic priority, a search of the works citing
\cite{KumariAdhikari2022} and of the literature on the negative spectrum of
the partial transpose (Section~\ref{sec:related}) found neither the
criterion of Theorem~\ref{thm:criterion}, nor the counterexamples, nor the
family $\rho_t$. The pure-state counterexample is likely familiar to anyone
who has computed both quantities. The search did not cover the literature on
rotationally invariant states, where $Q^\G$ and its spectrum probably appear
in another form. This is a statement about the documented search only, not a
claim of priority.

\clearpage
\begingroup\small
\begin{thebibliography}{99}
\setlength{\itemsep}{2pt}
\bibitem{KumariAdhikari2022}
A. Kumari and S. Adhikari,
\emph{Structured negativity: A physically realizable measure of entanglement
based on structural physical approximation},
Ann. Phys. \textbf{446}, 169113 (2022);
\href{https://doi.org/10.1016/j.aop.2022.169113}{doi:10.1016/j.aop.2022.169113};
\href{https://arxiv.org/abs/2209.03909}{arXiv:2209.03909}.

\bibitem{Peres1996}
A. Peres,
\emph{Separability criterion for density matrices},
Phys. Rev. Lett. \textbf{77}, 1413 (1996);
\href{https://doi.org/10.1103/PhysRevLett.77.1413}{doi:10.1103/PhysRevLett.77.1413};
\href{https://arxiv.org/abs/quant-ph/9604005}{arXiv:quant-ph/9604005}.

\bibitem{Horodecki1996}
M. Horodecki, P. Horodecki and R. Horodecki,
\emph{Separability of mixed states: necessary and sufficient conditions},
Phys. Lett. A \textbf{223}, 1 (1996);
\href{https://doi.org/10.1016/S0375-9601(96)00706-2}{doi:10.1016/S0375-9601(96)00706-2};
\href{https://arxiv.org/abs/quant-ph/9605038}{arXiv:quant-ph/9605038}.

\bibitem{VidalWerner2002}
G. Vidal and R. F. Werner,
\emph{A computable measure of entanglement},
Phys. Rev. A \textbf{65}, 032314 (2002);
\href{https://doi.org/10.1103/PhysRevA.65.032314}{doi:10.1103/PhysRevA.65.032314};
\href{https://arxiv.org/abs/quant-ph/0102117}{arXiv:quant-ph/0102117}.

\bibitem{HorodeckiEkert2002}
P. Horodecki and A. Ekert,
\emph{Direct detection of quantum entanglement},
Phys. Rev. Lett. \textbf{89}, 127902 (2002);
\href{https://doi.org/10.1103/PhysRevLett.89.127902}{doi:10.1103/PhysRevLett.89.127902};
\href{https://arxiv.org/abs/quant-ph/0111064}{arXiv:quant-ph/0111064}.

\bibitem{Kumari2023}
A. Kumari,
\emph{Detection and classification of bipartite and multipartite entangled
states}, doctoral thesis (2023);
\href{https://arxiv.org/abs/2305.16643}{arXiv:2305.16643}.

\bibitem{Aggarwal2023}
S. Aggarwal, A. Kumari and S. Adhikari,
\emph{Physical realization of realignment criteria using structural physical
approximation},
Phys. Rev. A \textbf{108}, 012422 (2023);
\href{https://doi.org/10.1103/PhysRevA.108.012422}{doi:10.1103/PhysRevA.108.012422};
\href{https://arxiv.org/abs/2301.09884}{arXiv:2301.09884}.

\bibitem{Rana2013}
S. Rana,
\emph{Negative eigenvalues of partial transposition of arbitrary bipartite
states},
Phys. Rev. A \textbf{87}, 054301 (2013);
\href{https://doi.org/10.1103/PhysRevA.87.054301}{doi:10.1103/PhysRevA.87.054301};
\href{https://arxiv.org/abs/1304.6775}{arXiv:1304.6775}.

\bibitem{Johnston2013}
N. Johnston,
\emph{Non-positive-partial-transpose subspaces can be as large as any
entangled subspace},
Phys. Rev. A \textbf{87}, 064302 (2013);
\href{https://doi.org/10.1103/PhysRevA.87.064302}{doi:10.1103/PhysRevA.87.064302};
\href{https://arxiv.org/abs/1305.0257}{arXiv:1305.0257}.

\bibitem{JiangQiuYu2026}
W.-J. Jiang, J.-T. Qiu and X.-D. Yu,
\emph{Spectral bounds for the partial transpose},
preprint (2026);
\href{https://arxiv.org/abs/2609.06617}{arXiv:2609.06617}.
\end{thebibliography}
\endgroup
\end{document}
